\spread{Can a complete human answer teach the reader where to look?}
\noindent A reference answer often uses words that do not appear together in the story. We trained the reader on the probability of producing that entire answer, including its final stop.
\platehere{Figure_182_learning_from_whole_replies}{Two questions about one previously unseen story section}{The question and human reference are followed by the original output pixels from the parent, the single-path reader and the marginal reader. Each reply was drawn from pictured inputs. We selected the first stopped pair that fit the display and changed the parent's semantic credit; if no credit changed, we used the first change in output.}
\begin{multicols}{2}
\section*{Learning without marked spans}
The frozen reader proposed up to 32 locations in a pictured page. Let $p_s$ be the learned probability of location $s$. For each location we evaluated $L_s$, the probability that the writer emits any of the distinct human reference images, followed by a stop. One head minimized $-\log\max_s p_s L_s$. A second minimized $-\log\sum_s p_s L_s$. Both models picked their highest-scoring location and generated a single reply at test time. No human wording entered this decision.
\section*{A matched experiment}
The two heads saw the same 10,084 questions in the same order during each of eight passes. They differed only in how the training objective combined locations. There were 5,988 narrative questions from 184 stories, 2,048 English statements and 2,048 Chinese exam questions. Computing 401,632 strip--wording probabilities took 180 seconds and 29.34 kJ on one RTX 4090 D. Each additional selector fit used less than 0.4 kJ. Earlier training, host power and assessment are outside these measurements.
\section*{What the reader wrote}
Asked who the bonnie lassie was, the single-path model chose a girdleful of cakes. The marginal model recovered the phrase “her only daughter” inside a grammatically defective sentence. The automated judge accepted the latter in both orders. The paired question asked how the widow prepared to bake. All three models echoed the question and omitted the bakeboard and meal.
\section*{What was held out}
We reserved 48 entire stories before fitting. Each contributed two different questions about the same section. All 96 questions and all 48 pictured sections were absent from the fit. Fourteen questions had at least one complete reference wording elsewhere in the training material. The held questions represent narrative reading of English in aligned glyph cells. We measured Chinese retention on a previously reserved examination set.
\end{multicols}
\newpage
\spread{When does a better probability become a better answer?}
\noindent The training objective scores complete human replies across the distribution of possible evidence. At deployment the reader chooses one strip, and the writer chooses one glyph at a time. We measured both the probabilities and the replies these decisions produced.
\platehere{Figure_183_full_answer_learning}{Training probabilities and autonomous meaning}{\textbf{a}, Running mean full-answer negative log likelihood during each pass through the 10,084 training questions. Both objectives use identical exposures. \textbf{b}, Stopped answers credited in both presentation orders by the same automated Qwen3 8B judge. Each group contains 16 questions. Qwen3-VL 8B Q4 received the same pictured content.}
\begin{multicols}{2}
\section*{Probability on unseen stories}
Across all 96 questions the mean negative log probability of a complete accepted answer was 70.504 nats for the parent, 70.132 for the single path and 70.104 for the marginal fit. The paired reduction against the parent was 0.400 nats (95\% story-bootstrap interval 0.157 to 0.636). Against the matched fit it was 0.029 nats ($-0.005$ to 0.063). This comparison supplies the human answer during probability measurement. In free generation, the parent and the two new systems each wrote zero verbatim answers on the 96 questions. The parent stopped on 92, the single-path model on 90 and the marginal model on 86. Stopped replies contained a complete reference wording on eight, eight and nine questions, respectively.
\section*{Meaning and transfer}
On the first 16 new story questions, the models earned one, one and two semantic credits. The quantized vision comparator earned ten. Across all 64 assessed questions the counts were 14, 11, 13 and 41. The automated judge accepted at least 15 of 16 human references and rejected all sibling foils in both orders in every group. The Chinese examination set yielded 58, 57 and 56 complete exact answers, respectively, out of 192. Chinese MLQA gave 20, 16 and 16. All three produced zero exact long-form answers in 128 trials and zero correct historical specimens in 260 period requests.
\section*{Which part must change?}
A Lean-checked identity separates two counts: questions for which every candidate leads the frozen writer to an unacceptable reply, and questions on which another candidate could succeed. On the earlier Chinese examination set, these account for 53 and 83 of the competing model's 136 exact-answer errors. On 92 earlier story questions the corresponding counts are 92 and zero. This partition depends on the chosen acceptance test; a faithful paraphrase can be correct. It identifies a necessary intervention for exact narrative composition under this interface: change the evidence available to the writer, its parameters or its decoder.
\end{multicols}
\newpage
